\chapter{시각화 --- 요약이 숨긴 것을 보기}
\label{ch:visualization}
% source: deck 5 slide 1

\begin{inbrief}
  \item[Why] 평균이나 표준편차 같은 \term{요약통계}(Summary Statistics) 하나만으로는
             분포의 \emph{모양}(봉우리(Peak) 개수, 치우침(Skewness),
             이상치(Outlier))이 지워진다. 시각화는 그 지워진 정보를 되찾는 도구다.
  \item[When] 변수 하나의 분포를 처음 확인할 때, 또는 여러 집단·변수의 분포를
              \emph{비교}할 때.
  \item[Where] 이 책 전체의 탐색 단계에 쓰인다. 다만 형태마다 전제가 다르다
               (예: 박스플롯(Box Plot)은 이봉분포(Bimodal Distribution)에 쓰면
               안 된다).
  \item[How] 먼저 ``이 그림으로 무엇을 묻는가''를 정한다 $\to$ 그 질문에 맞는
             형태를 고른다(\Cref{tab:form-choice}) $\to$ 형태마다의 함정
             (구간 개수, 에러바(Error Bar)가 무엇을 뜻하는지, 표본 크기)을
             확인한다.
\end{inbrief}

\Cref{ch:descriptive}은 \analogy{지도}가 반드시 무언가를 버린다는 결론으로
끝났다. 이 장은 그 버려진 것을 되찾는 \analogy{항공사진}을 다룬다.

강의의 표현을 그대로 옮기면 이렇다. \emph{``통계값은 숫자이고 유용하지만, 이야기
전체를 들려주지는 않는다.''} 시각화는 한 변수의 분포를 이해할 때 유용하고,
여러 변수의 분포를 \textbf{비교}할 때는 거의 필수다.

\begin{keyidea}
차트는 장식이 아니라 \textbf{질문에 답하는 도구}다. 먼저 ``이 그림으로 무엇을
묻는가''를 정하고, 그 다음에 형태를 고른다. 순서를 뒤집으면 예쁘지만 아무것도
답하지 않는 그림이 나온다.
\end{keyidea}

% =====================================================================
\section{배경}
\label{sec:vis-background}
% source: deck 5 "Basic Statistics and Visualization" slide 1

이 장은 강의 노트 \texttt{5\_StatandVisualization.pptx}(``Basic Statistics and
Visualization'')를 옮긴다. 이 장을 읽는 데 필요한 개념을 먼저 깐다. 이미
다룬 것은 한 줄 요약과 링크로, 이 장 안에서 처음 정의하는 것은 뒤의 각 절에서
다룬다.

\begin{itemize}
  \item \term{평균}(Mean)과 \term{표준편차}(Standard Deviation, 퍼짐의 크기)는
        \Cref{sec:center,sec:spread}에서 다뤘다. \Cref{sec:ci}의 신뢰구간
        공식 $\bar{x} \pm z \cdot s/\sqrt{n}$이 이 둘 위에 선다.
  \item \term{확률밀도함수}(Probability Density Function, PDF)는
        \Cref{sec:pmf-pdf}에서 다뤘다. \Cref{sec:histogram}의 히스토그램은
        이 PDF를 데이터로부터 경험적으로 근사한 것이다.
  \item \term{정규분포}(Normal Distribution)와 그 표준편차 배수(예:
        $z=1.96$)는 \Cref{sec:normal}에서 다룬다. \Cref{sec:ci}의 공식에
        쓰이는 $z$ 값이 여기서 나온다.
\end{itemize}

% =====================================================================
\section{분위수와 백분위수}
\label{sec:quantile}

\begin{inbrief}
  \item[Why] 평균 하나로는 부족하고, 분포의 \emph{여러 지점}을 알아야 하기 때문이다.
  \item[When] 박스플롯, $\IQR$, 신뢰구간 --- 이 장의 거의 모든 도구가 분위수 위에 선다.
  \item[Where] 왜도가 있는 데이터에서 특히 중요하다. 평균은 무너져도 분위수는 버틴다.
  \item[How] 정렬 $\rightarrow$ 각 값에 $0$에서 $1$까지의 위치를 매김 $\rightarrow$
             원하는 위치에서 보간(Interpolation).
\end{inbrief}

\begin{definition}[분위수와 백분위수]
\label{def:quantile}
\term{분위수}(Quantile)는 데이터를 \textbf{같은 크기의 그룹으로 나누는 경계선}이다.
\term{백분위수}(Percentile)는 그중 데이터를 \textbf{100개} 그룹으로 나누는 분위수다.

$k$번째 백분위수는 값들을 나누어, $k\%$ 가 그 아래에 놓이고 $(100-k)\%$ 가 그
위에 놓이게 하는 값이다.
\end{definition}

\begin{aside}[둘의 관계]
백분위수는 분위수의 \textbf{부분집합}이다. 분위수 \contains{} 백분위수.
같은 지점을 부르는 이름이 둘인 셈이라 혼용되지만, 엄밀히는 \emph{눈금의 개수}가
다르다 --- 사분위수(Quartile)는 4등분, 십분위수(Decile)는 10등분, 백분위수는
100등분이다. 소프트웨어에서는 보통 $0 \sim 1$ 을 받는 \texttt{quantile()} 하나로
통일되어 있다.
\end{aside}

\begin{center}
\begin{tabularx}{\linewidth}{@{}l l l X@{}}
\toprule
이름 & 몇 등분 & 표기 & 비고 \\
\midrule
사분위수(Quartile) & 4 & $Q_1, Q_2, Q_3$ & $Q_1$ = 하사분위수(Lower Quartile) \\
십분위수(Decile) & 10 & $D_1 \dots D_9$ & \\
백분위수(Percentile) & 100 & $P_1 \dots P_{99}$ & $P_{50} = Q_2 = $ 중앙값 \\
\bottomrule
\end{tabularx}
\end{center}

\begin{pitfall}[0\% 분위수는 왜 최솟값인가]
강의가 던지는 질문이다. $0\%$ 분위수(0번째 백분위수)는 \textbf{최솟값}이다.
정의상 ``그 값보다 작은 값이 $0\%$''여야 하는데, 최솟값보다 작은 값은 데이터에
하나도 없기 때문이다. 같은 이유로 $100\%$ 분위수는 최댓값이다.
\end{pitfall}

\subsection{직접 계산해 보기}
\label{subsec:quantile-calc}

\begin{example}[25\% 분위수 구하기]
\label{ex:quantile}
강의 예제다. 여섯 개의 값이 있다: $3.7,\ 2.7,\ 3.3,\ 1.3,\ 2.2,\ 3.1$.

\textbf{1단계 --- 정렬한다.}
\[
  1.3,\quad 2.2,\quad 2.7,\quad 3.1,\quad 3.3,\quad 3.7
\]

\textbf{2단계 --- 각 값에 $0$부터 $1$까지 등간격의 위치를 매긴다.}
$n = 6$ 이므로 $i$번째 값의 위치는 $(i-1)/(n-1) = (i-1)/5$ 다.

\begin{center}
\begin{tabular}{@{}lrrrrrr@{}}
\toprule
$i$ & 1 & 2 & 3 & 4 & 5 & 6 \\
값 & 1.3 & 2.2 & 2.7 & 3.1 & 3.3 & 3.7 \\
위치 & 0.0 & 0.2 & 0.4 & 0.6 & 0.8 & 1.0 \\
\bottomrule
\end{tabular}
\end{center}

\textbf{3단계 --- 0.25는 눈금 사이에 있으므로 보간한다.}
$0.25$ 는 $0.2$(값 2.2)와 $0.4$(값 2.7) 사이에 있다. $0.2$ 에서
$0.05$ 만큼 갔으므로 전체 간격 $0.2$ 의 $1/4$ 지점이다.
\[
  0.75 \times 2.2 + 0.25 \times 2.7 = 1.65 + 0.675 = \mathbf{2.325}
\]

즉 하사분위수는 2.325다. \textbf{데이터에 없는 값}이라는 점에 주의하자 ---
분위수는 관측값 중 하나일 필요가 없다.
\end{example}

\begin{lstlisting}[language=Python]
import pandas as pd
s = pd.Series([3.3, 2.7, 2.2, 1.3, 3.7, 3.1])
s.quantile(0.25)    # 2.325  <- 위 손계산과 일치
\end{lstlisting}

\begin{figure}[ht]\centering
  \includegraphics[width=0.85\linewidth]{05_01_quantiles.png}
  \caption{정렬된 값에 위치를 매기고 원하는 지점을 읽는 과정. 곡선의 가로축이
  \Cref{ex:quantile}의 ``위치'', 세로축이 그 위치의 값이다.}
  \label{fig:quantile-curve}
\end{figure}

% =====================================================================
\section{히스토그램}
\label{sec:histogram}

\begin{inbrief}
  \item[Why] 확률밀도함수(Probability Density Function)를 데이터로부터 경험적으로
             그려 보기 위해서다.
  \item[When] 변수 하나의 분포 모양을 처음 확인할 때. 가장 먼저 그리는 그림이다.
  \item[Where] 연속형 변수에 쓴다. 범주형에는 막대그래프(Bar Chart)를 쓴다.
  \item[How] 데이터를 구간(Bin)으로 나누고, x축에 구간 중심을, y축에 그 구간의
             개수(또는 비율)를 놓는다.
\end{inbrief}

\term{히스토그램}(Histogram)은 확률밀도함수의 \textbf{경험적 표현}(Empirical
Representation)이다. 만드는 법은 단순하다.

\begin{enumerate}
  \item 데이터를 구간(Bin)으로 나눈다.
  \item x축 = 구간의 중심점.
  \item y축 = 그 구간에 들어간 데이터 개수(또는 전체 대비 비율).
\end{enumerate}

\subsection{구간 개수가 결론을 바꾼다}
\label{subsec:binning}

강의는 ``구간을 어떻게 나눌 것인가?''를 괄호 안에 던져 놓고 지나간다. 그런데
이것이 히스토그램의 유일한 자유도이자 가장 큰 함정이다.

\begin{figure}[ht]\centering
  \includegraphics[width=\linewidth]{05_02_histogram_bins.png}
  \caption{같은 데이터, 다른 구간 개수. 구간이 너무 적으면 구조가 뭉개지고, 너무
  많으면 표본 잡음(Sampling Noise)이 봉우리처럼 보인다.}
  \label{fig:bins}
\end{figure}

\begin{pitfall}[구간 개수는 보고해야 할 선택이다]
구간 개수를 바꾸면 봉우리가 생겼다 사라진다. 즉 \textbf{같은 데이터에서 다른
결론}이 나온다. 따라서 히스토그램에서 읽은 구조(``봉우리가 두 개다'')를
주장할 때는 다른 구간 개수에서도 그 구조가 살아남는지 확인해야 한다.
\Cref{fig:two-populations}의 두 집단은 구간을 어떻게 잡아도 사라지지 않기
때문에 믿을 수 있는 구조다.
\end{pitfall}

\subsection{여러 분포를 나란히 놓기}
\label{subsec:small-multiples}

히스토그램 여러 개를 같은 축으로 나란히 배치하면 집단 간 차이가 드러난다. 강의는
서브레딧별 순위 분포를 예로 든다.

\begin{figure}[ht]\centering
  \includegraphics[width=\linewidth]{05_03_small_multiples.png}
  \caption{작은 배수(Small Multiples). 축을 공유해야 비교가 성립한다 --- 패널마다
  축 범위가 다르면 눈은 그 차이를 읽지 못한다.}
  \label{fig:small-multiples}
\end{figure}

\begin{pitfall}[표본 크기를 같이 적는다]
\Cref{fig:small-multiples}에서 버거는 69개, 나머지는 11--12개다. 개수가 적은
패널은 모양이 들쭉날쭉한데, 이것은 분포의 성질이 아니라 \textbf{표본이 작아서}
생긴 것이다. 각 패널에 $n$ 을 적어 두지 않으면 독자가 잡음을 구조로 읽는다.
\end{pitfall}

% =====================================================================
\section{선 그래프 --- 밀도에 더 가까이}
\label{sec:density-line}

\begin{inbrief}
  \item[Why] 히스토그램의 계단을 매끄럽게 만들어 확률밀도함수에 더 가까이 가기 위해서다.
  \item[When] 두 개 이상의 분포를 \emph{겹쳐서} 비교할 때. 히스토그램은 겹치면 가린다.
  \item[Where] 데이터가 충분히 많을 때. 표본이 적으면 매끄러운 곡선이 없는 구조를
               지어낸다.
  \item[How] 구간을 충분히 잘게 나눠 곡선이 매끄러우면서도 분포를 왜곡하지 않는
             지점을 찾는다.
\end{inbrief}

강의의 기준은 명확하다. \emph{``곡선이 충분히 매끄러우면서도 분포를 정확히
표현하도록, 충분한 수의 구간을 쓴다.''} 두 조건이 서로 반대 방향으로 당기므로
절충점을 찾아야 한다 --- \Cref{subsec:binning}의 문제가 그대로 이어진다.

\begin{figure}[ht]\centering
  \includegraphics[width=\linewidth]{05_04_density_line.png}
  \caption{밀도 곡선(Density Line). 히스토그램의 계단을 매끄럽게 편 것이며,
  여러 분포를 겹쳐 비교할 때 서로를 가리지 않는다는 장점이 있다.}
  \label{fig:density}
\end{figure}

% =====================================================================
\section{에러바 --- 요약을 나란히 놓기}
\label{sec:errorbar}

\begin{inbrief}
  \item[Why] 비교할 범주가 너무 많아 분포를 전부 그릴 수 없을 때, 각 분포의
             \emph{요약}만 나란히 놓기 위해서다.
  \item[When] 연속 변수를 구간으로 나눠 그 함수로 볼 때, 또는 집단이 여럿일 때.
  \item[Where] 분포가 단봉이고 대칭일 때만 정직하다. 이봉분포는 요약이 거짓말을 한다.
  \item[How] 중심(평균)을 점이나 선으로, 퍼짐을 막대로 그린다.
\end{inbrief}

\begin{figure}[ht]\centering
  \includegraphics[width=\linewidth]{05_05_error_bar.png}
  \caption{집단별 평균 곡선과 퍼짐을 나타내는 막대. 강의의 예는 약물에 저항성인
  환자군과 민감한 환자군의 비교다.}
  \label{fig:errorbar}
\end{figure}

\begin{figure}[ht]\centering
  \includegraphics[width=\linewidth]{05_06_error_variants.png}
  \caption{에러바의 변형. 왼쪽은 에러 구름(Error Cloud), 오른쪽은 에러바를 얹은
  막대그래프(Bar Plot with Error Bars)다. 담는 정보는 같고 강조점이 다르다.}
  \label{fig:error-variants}
\end{figure}

\begin{pitfall}[막대가 무엇인지 캡션에 반드시 쓴다]
\label{pit:sd-se-ci}
에러바의 길이로 쓰이는 후보가 셋이고, \textbf{길이가 서로 몇 배씩 차이 난다}.

\begin{center}
\begin{tabularx}{\linewidth}{@{}l l X@{}}
\toprule
무엇 & 계산 & 무엇을 말하는가 \\
\midrule
표준편차 $\sd$ & $\sigma$ & \textbf{데이터}가 얼마나 퍼져 있는가 \\
표준오차 $\mathrm{SE}$ & $\sigma/\sqrt{n}$ & \textbf{평균 추정값}이 얼마나 불안정한가 \\
신뢰구간 $\mathrm{CI}$ & $\approx \bar{x} \pm 1.96\,\mathrm{SE}$ & 참평균이 있을 법한 범위 \\
\bottomrule
\end{tabularx}
\end{center}

버거 데이터로 보면 $\sd = 277.16$ 이지만 $\mathrm{SE} = 277.16/\sqrt{69} = 33.37$
로 \textbf{8배 넘게 차이} 난다. 캡션에 무엇인지 쓰지 않은 에러바는 읽을 수 없다.
\end{pitfall}

% =====================================================================
\section{바이올린 플롯}
\label{sec:violin}

\begin{inbrief}
  \item[Why] 에러바와 박스플롯이 버리는 \emph{분포의 모양}을 되살리기 위해서다.
  \item[When] 집단을 비교하는데 그 분포가 단봉이 아닐 가능성이 있을 때.
  \item[Where] 표본이 충분해야 한다. $n$ 이 작으면 밀도 추정이 없는 봉우리를 만든다.
  \item[How] 각 집단의 밀도 곡선을 세로로 세우고 좌우 대칭으로 펼친다.
\end{inbrief}

\begin{figure}[ht]\centering
  \includegraphics[width=\linewidth]{05_07_violin.png}
  \caption{바이올린 플롯. 폭이 곧 그 값 근처의 밀도이므로, 봉우리가 몇 개인지가
  즉시 보인다. 강의의 예는 보스턴 마라톤 완주 시간의 연도별 분포다.}
  \label{fig:violin}
\end{figure}

\begin{figure}[ht]\centering
  \includegraphics[width=\linewidth]{05_08_violin_combined.png}
  \caption{결합 바이올린 플롯. 두 집단을 좌우 반쪽씩 맞붙여 같은 축에서 비교한다.
  두 집단을 나란히 두 개 그리는 것보다 시선 이동이 짧다.}
  \label{fig:violin-combined}
\end{figure}

% =====================================================================
\section{바코드 차트}
\label{sec:barcode}

\begin{inbrief}
  \item[Why] 요약하지 않고 \textbf{모든 데이터 점}을 그대로 보여 주기 위해서다.
  \item[When] 데이터 개수가 적어 요약이 오히려 정보를 없앨 때.
  \item[Where] 각 행의 값 범위가 서로 비교 가능해야 한다. 아니면 행끼리 읽을 수 없다.
  \item[How] 데이터 점 하나를 선 위의 막대 하나로 그리고, 변수마다 행을 따로 준다.
\end{inbrief}

강의의 예는 워싱턴 메트로의 정시 운행률이다. 각 행이 한 노선, 각 행의 12개 막대가
월별 평균이고, 연평균은 굵은 막대로 표시된다.

\begin{figure}[ht]\centering
  \includegraphics[width=\linewidth]{05_09_barcode.png}
  \caption{바코드 차트. 점 하나하나가 보이므로 \emph{빈 구간}이 눈에 띈다 ---
  요약통계로는 절대 드러나지 않는 정보다.}
  \label{fig:barcode}
\end{figure}

\begin{keyidea}
바코드 차트는 이 장에서 유일하게 \textbf{아무것도 버리지 않는} 그림이다.
\analogy{지도}가 아니라 \analogy{지형 그 자체}다. 그래서 데이터가 조금만 많아져도
쓸 수 없지만, 적을 때는 가장 정직하다.
\end{keyidea}

% =====================================================================
\section{박스플롯으로 집단 비교하기}
\label{sec:boxplot}

\begin{inbrief}
  \item[Why] 여러 집단의 중심과 퍼짐을 한 화면에서 비교하기 위해서다.
  \item[When] 집단이 3개 이상이고 각 집단의 분포가 대체로 단봉일 때.
  \item[Where] 이봉분포에는 쓰면 안 된다 --- 박스플롯은 봉우리 개수를 표현하지 못한다.
  \item[How] $Q_1$ 에서 $Q_3$ 까지 상자를 그리고, $1.5 \times \IQR$ 이내의 가장 먼
             관측값까지 수염을 뻗고, 그 밖은 이상치로 찍는다.
\end{inbrief}

박스플롯을 만드는 규칙은 \Cref{subsec:fence}에서 이미 다뤘다. 여기서는 그것을
\emph{비교}에 쓰는 법을 본다.

\begin{figure}[ht]\centering
  \includegraphics[width=\linewidth]{05_11_box_compare.png}
  \caption{박스플롯으로 집단 비교하기. 중앙선의 높이 차이가 중심의 차이를,
  상자의 길이 차이가 퍼짐의 차이를 나타낸다.}
  \label{fig:box-compare}
\end{figure}

강의는 붓꽃 데이터에서 두 가지를 읽는다 --- 꽃받침(Sepal)의 중앙값이 꽃잎(Petal)
보다 높고, 꽃잎 쪽의 값의 퍼짐이 더 크다. 이 두 문장이 박스플롯에서 읽어야 할
것의 전부다. \textbf{중심의 차이}와 \textbf{퍼짐의 차이}.

과제 데이터에서도 같은 방식으로 읽힌다. 음식 유형별 칼로리 중앙값은
치킨너겟 240 $\rightarrow$ 감자튀김 315 $\rightarrow$ 그릴드치킨샌드위치 398
$\rightarrow$ 브레디드치킨샌드위치 510 $\rightarrow$ 밀크셰이크 585
$\rightarrow$ 버거 591 로 단조 증가한다.

\subsection{노치 박스플롯}
\label{subsec:notch}

\begin{figure}[ht]\centering
  \includegraphics[width=\linewidth]{05_12_box_notched.png}
  \caption{노치(Notch)를 넣은 박스플롯. 허리가 잘록한 부분이 중앙값의 신뢰구간을
  나타내며, 두 상자의 노치가 겹치지 않으면 중앙값이 다르다고 볼 근거가 된다.}
  \label{fig:box-notch}
\end{figure}

\term{노치}(Notch)는 중앙값 주변의 불확실성을 상자에 직접 새긴 것이다. 두 집단의
노치가 겹치지 않으면 중앙값이 통계적으로 다르다고 본다 --- \Cref{sec:ci}의
신뢰구간 비교를 박스플롯 안으로 가져온 것이다.

\begin{pitfall}[노치가 상자보다 커지면 표본이 부족한 것]
노치의 폭은 $\pm 1.58 \times \IQR/\sqrt{n}$ 으로 계산된다. $n$ 이 작으면 이 값이
$\IQR$ 자체보다 커져서 \textbf{노치가 상자 밖으로 뒤집혀 나온다}. 그림이
망가진 것이 아니라, ``이 표본으로는 중앙값을 말할 수 없다''는 경고다.
\end{pitfall}

% =====================================================================
\section{신뢰구간}
\label{sec:ci}

\begin{inbrief}
  \item[Why] 계산한 통계값에 \emph{얼마나 불확실성이 있는지}를 함께 말하기 위해서다.
  \item[When] 표본에서 계산한 값으로 모집단을 말하려 할 때 --- 즉 거의 항상.
  \item[Where] 두 집단이 ``정말로'' 다른지 판단할 때. 겹치는지만 보면 된다.
  \item[How] 공식($\bar{x} \pm z \cdot \sigma/\sqrt{n}$) 또는 부트스트랩
             (Bootstrap) 둘 중 하나로 구한다.
\end{inbrief}

\term{신뢰구간}(Confidence Interval)은 특정 통계값에 얼마나 불확실성이 있는지를
나타낸다. 설문 조사에서 ``오차 범위''(Margin of Error)라고 부르는 것과 같다 ---
전체 모집단을 조사했다면 얻었을 값이 이 범위 안에 있으리라는 뜻이다.

\subsection{공식으로 구하기}

\[
  \mathrm{CI} = \bar{x} \pm z \cdot \frac{s}{\sqrt{n}}
\]

\begin{center}
\begin{tabularx}{\linewidth}{@{}l X@{}}
\toprule
기호 & 의미 \\
\midrule
$s$ & 표준편차(Standard Deviation) \\
$n$ & 관측값의 개수 \\
$z$ & 표에서 읽는 $z$ 값. $95\%$ 신뢰구간이면 $1.96$ \\
$\bar{x}$ & 신뢰구간을 구하려는 대상 값 \\
\bottomrule
\end{tabularx}
\end{center}

$z = 1.96$ 의 의미는 이렇다. 평균에서 $-1.96$ 부터 $+1.96$ 표준편차 사이의 값이
무작위로 뽑힐 확률이 $95\%$ 다. 뒤집으면, 어떤 값이 무작위로 뽑힐 확률이 $5\%$
미만이면 그것을 \textbf{통계적으로 유의하다}(Statistically Significant)고 본다.

\subsection{부트스트랩으로 구하기}
\label{subsec:bootstrap}

공식은 분포에 대한 가정을 필요로 한다. \term{부트스트랩}(Bootstrapping)은 그
가정 없이, 컴퓨터로 재표집을 반복해서 같은 것을 구한다. 강의의 예는 암컷 생쥐
12마리의 몸무게다.

\begin{figure}[ht]\centering
\begin{tikzpicture}[
  node distance = 6mm,
  bsbox/.style = {draw = bookrule, fill = booksurface, rounded corners = 2pt,
                  minimum width = 62mm, minimum height = 8.5mm, align = center,
                  font = \small},
  arrow/.style = {-{Stealth[length=2.2mm]}, draw = bookmuted}]
  \node[bsbox] (a) {표본 평균을 계산한다};
  \node[bsbox, below = of a] (b)
        {12개를 \textbf{복원추출}(Replacement)로 다시 뽑는다\\[-1pt]
         \footnotesize 같은 개체가 여러 번 뽑혀도 된다};
  \node[bsbox, below = of b] (c) {재표집된 12개의 평균을 계산한다};
  \node[bsbox, below = of c] (d) {2--3 단계를 \textbf{10{,}000회 이상} 반복한다};
  \node[bsbox, below = of d, fill = bookaccentbg, draw = bookaccent] (e)
        {10{,}000개 평균의 2.5\%\,--\,97.5\% 구간\\[-1pt]
         \footnotesize 이것이 95\% 신뢰구간이다};
  \draw[arrow] (a) -- (b);
  \draw[arrow] (b) -- (c);
  \draw[arrow] (c) -- (d);
  \draw[arrow] (d) -- (e);
  \draw[arrow] (d.east) -- ++(9mm,0) |- (b.east);
\end{tikzpicture}
\caption{부트스트랩 절차. 핵심은 \textbf{복원추출}이다 --- 중복을 허용하지 않으면
매번 같은 표본이 나와 아무 변동도 생기지 않는다.}
\label{fig:bootstrap-flow}
\end{figure}

\begin{figure}[ht]\centering
  \includegraphics[width=\linewidth]{05_13_bootstrap_ci.png}
  \caption{버거 칼로리의 부트스트랩 신뢰구간. 재표집 10{,}000회의 평균 분포이며,
  2.5 백분위수와 97.5 백분위수가 신뢰구간의 양끝이 된다.}
  \label{fig:bootstrap}
\end{figure}

\begin{pitfall}[이 히스토그램은 버거의 분포가 아니다]
\Cref{fig:bootstrap}의 히스토그램은 \textbf{10{,}000개 ``평균''의 분포}이지
69개 버거의 분포가 아니다. 가로축 범위가 555--690 정도로 좁은 이유가 여기 있다 ---
개별 버거는 140부터 1240까지 퍼져 있지만, 69개의 평균은 그만큼 흔들리지 않는다.
이 둘을 혼동하면 ``버거는 전부 600칼로리 근처''라는 완전히 틀린 결론에 이른다.
\end{pitfall}

\begin{keyidea}
표본 크기를 키우면 신뢰구간은 \textbf{좁아지고}, 표준편차가 크면 신뢰구간은
\textbf{넓어진다}. 공식 $\bar{x} \pm z \cdot s/\sqrt{n}$ 에서 $n$ 은 분모에,
$s$ 는 분자에 있기 때문이다.
\end{keyidea}

\subsection{두 집단을 비교하기}
\label{subsec:ci-compare}

\begin{figure}[ht]\centering
  \includegraphics[width=\linewidth]{05_14_ci_compare.png}
  \caption{신뢰구간으로 두 집단 비교하기. 두 구간이 겹치지 않으므로 두 평균은
  통계적으로 다르다($p < 0.05$).}
  \label{fig:ci-compare}
\end{figure}

버거의 평균 칼로리는 620이고 $95\%$ 신뢰구간은 $[557,\ 686]$, 치킨너겟은 평균
275에 신뢰구간 $[211,\ 345]$ 다. 두 구간 사이에는 345와 557 사이의 빈틈이 있다.
\textbf{겹치지 않으므로} 두 평균이 우연히 다르게 나왔을 가능성은 $5\%$ 미만이다.

\begin{pitfall}[겹친다고 같은 것은 아니다]
이 판정법은 한 방향으로만 강하다. 구간이 \textbf{겹치지 않으면} 다르다고 말할 수
있지만, \textbf{겹친다고 해서 같다}고 말할 수는 없다. 겹치는 경우는 ``이 데이터
로는 판단할 수 없다''는 뜻이고, 정확한 판정에는 별도의 검정이 필요하다.
두 평균을 직접 비교하는 검정은 \Cref{sec:cd-t-two}에서 다룬다.
\end{pitfall}

% =====================================================================
\section{형태 고르기 --- 정리}
\label{sec:choosing-form}

\begin{table}[ht]\centering
\caption{질문과 형태의 대응. 위에서 아래로 갈수록 데이터를 덜 버린다.}
\label{tab:form-choice}
\begin{tabularx}{\linewidth}{@{}l l X@{}}
\toprule
묻는 것 & 형태 & 그 형태인 이유 \\
\midrule
한 변수의 분포 & 히스토그램(Histogram) & 구간별 개수를 직접 센다 \\
분포를 겹쳐서 비교 & 밀도 곡선(Density Line) & 겹쳐도 서로를 가리지 않는다 \\
여러 집단의 분포 & 작은 배수(Small Multiples) & 축을 공유해 비교를 성립시킨다 \\
집단이 많을 때의 요약 & 에러바(Error Bar) & 중심과 퍼짐만 남겨 자리를 아낀다 \\
집단의 중심과 퍼짐 & 박스플롯(Box Plot) & 다섯 개 숫자로 요약해 나란히 놓는다 \\
중앙값의 차이 판정 & 노치 박스플롯(Notched Box Plot) & 불확실성을 상자에 새긴다 \\
분포의 \emph{모양}까지 & 바이올린(Violin Plot) & 박스가 버린 봉우리를 되살린다 \\
모든 점을 그대로 & 바코드(Barcode Chart) & 아무것도 버리지 않는다 \\
추정값의 불확실성 & 신뢰구간(Confidence Interval) & 겹치는지로 차이를 판정한다 \\
\bottomrule
\end{tabularx}
\end{table}

\section{요약}

\begin{itemize}
  \item 분위수는 정렬 $\rightarrow$ 위치 부여 $\rightarrow$ 보간으로 계산되며,
        관측값 중 하나일 필요가 없다(\Cref{ex:quantile}).
  \item 히스토그램의 구간 개수는 \emph{선택}이고, 선택이 결론을 바꾼다.
  \item 에러바의 막대가 $\sd$인지 $\mathrm{SE}$인지 $\mathrm{CI}$인지 반드시
        밝힌다. 길이가 몇 배씩 차이 난다(\Cref{pit:sd-se-ci}).
  \item 박스플롯은 봉우리 개수를 표현하지 못한다. 의심되면 바이올린이나 바코드로
        확인한다.
  \item 신뢰구간이 겹치지 않으면 다르다. 겹친다고 같은 것은 아니다.
  \item 부트스트랩은 분포 가정 없이 신뢰구간을 구한다. 핵심은 복원추출이다.
\end{itemize}

% =====================================================================
\section{어떻게 적용하는가}
\label{sec:vis-apply}

이 장의 차트와 개념은 이 책의 실전 장에서 되풀이해 쓰인다.

\begin{itemize}
  \item \Cref{ch:hw1}(과제 1 사례집)은 버거와 밀크셰이크의 칼로리
        분포를 이 장의 차트가 아니라 평균·SD·최소·최대로 된 요약통계 표로
        비교한다(\Cref{tab:hw1-burger-vs-shake}). 버거와 치킨너겟을 신뢰구간으로
        직접 겹쳐 보는 비교는 이 장의 \Cref{subsec:ci-compare}에서 한다.
  \item \Cref{ch:hw2}(과제 2, 분포 비교)는 영화 평점과 나이 분포에는 이 장의
        히스토그램을 쓰지만, 공항 노선 수는 꼬리가 두 자릿수 넘게 퍼져
        히스토그램 대신 로그-로그 축의 CCDF(Complementary Cumulative
        Distribution Function)를 쓴다.
  \item \Cref{sec:qq}(Q-Q 플롯, 강의 7)은 \Cref{ex:quantile}과 같은 분위수
        비교 아이디어를 쓰지만, 플로팅 위치(Plotting Position) 규칙은 다르다.
        \Cref{ex:quantile}은 $(i-1)/(n-1)$을, \Cref{sec:qq}는 $p_i = i/(n+1)$을
        쓴다(\Cref{sec:qq}의 함정에서 둘을 대조한다).
  \item \Cref{sec:cd-t-two}(강의 11)의 두 표본 $t$검정(Two-Sample $t$-Test)은
        이 장의 신뢰구간 겹침 판정을 평균 비교로 정식화한 것이다.
        \Cref{sec:cd-ks-two}의 두 표본 KS검정(Kolmogorov-Smirnov Test)은
        평균이 아니라 분포 모양 전체를 비교한다.
\end{itemize}

% =====================================================================
\section{앞으로}
\label{sec:vis-future}

\Cref{subsec:ci-compare}의 신뢰구간 판정과 $5\%$ 유의 기준은 이 장에서는
직관으로만 제시된다. 그 직관을 귀무가설(Null Hypothesis)·유의수준
(Significance Level)·기각 규칙(Rejection Rule)으로 정식화한 것이
\Cref{ch:hypothesis}다. 또한 \Cref{subsec:ci-compare}의 함정 상자가 던진
``구간이 겹치면 같은 것인가''라는 질문도 \Cref{ch:hypothesis}가 답한다.

신뢰구간 겹침 판정은 평균 하나만 비교한다. 평균만 더 정식으로 비교하려면
\Cref{sec:cd-t-two}의 두 표본 $t$검정(Two-Sample $t$-Test)이 필요하고,
분포 모양 전체를 비교하려면 \Cref{sec:cd-ks-two}의 두 표본
KS검정(Kolmogorov-Smirnov Test)이 필요하다.
